\section{立项依据}
\subsection{国内外研究现状}

目前，冷弯薄壁型钢结构分析主要采用杆系有限元模型和壳体有限元模型两类。杆系模型以梁单元（如 Euler-Bernoulli 梁、Timoshenko 梁）或杆单元模拟梁柱构件，计算效率高，适用于整体结构受力分析，但难以准确捕捉局部屈曲和畸变屈曲；壳体模型采用四边形壳单元（如 ABAQUS 的 S4R、ANSYS 的 SHELL181）对截面进行离散，能够精确描述局部屈曲、畸变屈曲及应力集中，但网格划分和计算成本显著增加。为兼顾精度与效率，国内外学者提出了广义梁理论（Generalized Beam Theory, GBT）、有限条法（Finite Strip Method, FSM）以及基于等几何分析（Isogeometric Analysis, IGA）的折中方案，并已在部分通用软件（如 ABAQUS、ANSYS、OpenSees）和专用程序中实现。

在实际分析中，冷弯薄壁型钢构件普遍存在以下问题：① \textbf{剪切自锁与膜自锁}：传统低阶单元在模拟薄壁梁时会出现剪切/膜锁定，需通过减缩积分、假设应变法或高阶单元加以消除；② \textbf{屈曲模式耦合}：局部屈曲、畸变屈曲与整体屈曲相互作用，需要合理的单元选择、缺陷引入和屈曲路径追踪；③ \textbf{初始缺陷与残余应力}：薄壁构件对初始几何缺陷、残余应力及边界条件极为敏感，需依据规范或实测数据合理取值；④ \textbf{连接半刚性}：螺钉、自攻螺钉等连接件具有显著的半刚性特征，节点域建模复杂。针对上述问题，现有研究已从理论模型、单元构造、本构关系、缺陷施加及连接模拟等方面提出了多种解决方案，为本项目的模块开发提供了理论基础。


\subsection{项目可行性分析}
\begin{itemize}[leftmargin=2em]
    \item \textbf{技术可行性：}团队中包括教授1名、副教授2名，长期致力于有限元领域、钢结构领域研究，具有丰富的有限元软件开发及结构分析经验。
    \item \textbf{软件可行性：}课题组具备自主开发程序 ApproxOperator.jl（\url{https://github.com/jmwjc/ApproxOperator.jl}），程序已具备框架结构静力、动力及屈曲分析能力，可作为本项目的算法验证与快速原型开发平台。
    \item \textbf{数据可行性：}合作企业大禾众邦可提供丰富的工程案例、真实截面参数、荷载工况及设计经验数据，为模块验证与规范校核提供数据保障。
    \item \textbf{硬件可行性：}实验室配备高性能计算服务器，可满足大规模有限元模型的计算需求；同时可依托 Julia、Python、Gmsh、ParaView 等开源科学计算生态，降低开发成本。
    \item \textbf{经济与时间可行性：}项目周期约为 8 个月，经费预算主要用于计算设备、学术交流、论文发表及知识产权申请，与项目目标相匹配，具备较强的可执行性。
\end{itemize}

